

%%%%%%%%%%%%%%%%%%%%%%%%%%%%%%%%%
%	 PACKAGES	%
%%%%%%%%%%%%%%%%%%%%%%%%%%%%%%%%%
% Codierung von in und Output
\usepackage[utf8]{inputenc}
\usepackage[T1]{fontenc}

% AMS-Pakete für Mathe
\usepackage{amsthm,amssymb,amsmath}

% Sprachpaket neue rechtschreibung
\usepackage[ngerman]{babel} 


\usepackage[many]{tcolorbox}

%Header und Seitengrösse
\usepackage{fancyhdr}
\usepackage{geometry}

% Paket für Links im PDF (glaubs)!
\usepackage[hidelinks]{hyperref}

%Paket um ganze pdf-File zu importieren (Arbeitsblätter aus Analysis DMK)
\usepackage{pdfpages}

%TColorbox nutzt Theoreme und bessere Verweise
\tcbuselibrary{theorems}
\usepackage{cleveref}
\crefname{section}{Kapitel}{Kapitel}
\Crefname{section}{Kapitel}{Kapitel}

%Package, um Aufgaben auf mehrere Spalten zu verteilen.
\usepackage{tasks}

\usepackage{xcolor}
\usepackage{listings}
\definecolor{codegreen}{rgb}{0,0.6,0}
\definecolor{codegray}{rgb}{0.5,0.5,0.5}
\definecolor{codepurple}{rgb}{0.58,0,0.82}
\definecolor{backcolour}{rgb}{0.95,0.95,0.92}
 
\lstdefinestyle{mystyle}{
    backgroundcolor=\color{backcolour},   
    commentstyle=\color{codegreen},
    keywordstyle=\color{magenta},
    numberstyle=\tiny\color{codegray},
    stringstyle=\color{codepurple},
    basicstyle=\footnotesize,
    breakatwhitespace=false,         
    breaklines=true,                 
    captionpos=b,                    
    keepspaces=true,                 
    numbers=left,                    
    numbersep=5pt,                  
    showspaces=false,                
    showstringspaces=false,
    showtabs=false,                  
    tabsize=2
}
\lstset{style=mystyle}

% kein Plan, was das folgende macht... :(
% \lstset{basicstyle=\ttfamily}
% \makeatletter
% \def\BState{\State\hskip-\ALG@thistlm}
% \makeatother

\usepackage{graphicx}
\usepackage{hyperref}
\usepackage{array}
\hypersetup{colorlinks=true}

% damit kann man nummerierte Liste unterbrechen und wieder aufnehmen
\usepackage{enumitem}
\newcounter{tempCounter}

% SI einheiten
\usepackage{siunitx}

%%%CHANGE LABEL FIGURE
\usepackage[figurename=Abbildung.]{caption}
\usepackage{subcaption}

%%%%%%%%%%%%%%%%%%%%%%%%%%%%%%%%
%%%%%%%%%%%	LAYOUT	%%%%%%%%%%%%
%%%%%%%%%%%%%%%%%%%%%%%%%%%%%%%%

\geometry{
	a4paper,
	total={165mm,225mm},
	left=25mm,
	top=35mm,
}

%keine Einrückung bei Absatz dafür Abstand
\setlength\parindent{0pt}
\setlength{\parskip}{1em}

%Header
\pagestyle{fancy}
\setlength{\headheight}{13.9pt}
\renewcommand{\sectionmark}[1]{ \markright{#1}{} }
\lhead[L]{M. Anderegg}
\chead{}
\rhead{\rightmark}
\lfoot{}
\cfoot{\thepage}
\rfoot{}

%Tabellen
\usepackage{booktabs}



%%%%%%%%%%%%%%%%%%%%%%%%%%%%%%%%
%%%%%%%%% TCB-Commands %%%%%%%%%
%%%%%%%%%%%%%%%%%%%%%%%%%%%%%%%%

\newtcbtheorem
[auto counter, number within=section, crefname={beispiel}{beispiele}, Crefname={Beispiel}{Beispiele}] % Nummerierung und Name bei Verweis
{bsp} % Name der Umgebung \begin{bsp}
{Beispiel} %Name der im Titel angezeigt wird
{breakable, enhanced, fonttitle=\bfseries, coltitle = black, colbacktitle = orange!15, righthand width=3cm,
	sharp corners,boxrule=.4pt,colback=orange!5} %%Optionen für die Box. 
	%Breakable und Enhanced lassen zu, dass sich die Box auf zwei Seiten verteilt.
{bsp} %So wird intern auf die Referenz verwiesen.

\newtcbtheorem[auto counter,  crefname={Einführungaufgaben}{Einführungaufgaben},
Crefname={Einführungaufgaben}{Einführungaufgaben}]{LA}{Einführungaufgaben}{
	breakable, enhanced, fonttitle=\bfseries, coltitle = black, colbacktitle = cyan!15, righthand width=3cm,
	sharp corners,boxrule=.4pt,colback=white}{LA}

\newtcbtheorem[auto counter, number within=section,  crefname={Übung}{Übungen}, Crefname={Übung}{Übungen}]{ueb}{Übung}{
	breakable, fonttitle=\bfseries, coltitle = black, colbacktitle = blue!15, righthand width=3cm,
	sharp corners,boxrule=.4pt,colback=blue!5}{uebng}


\newtcbtheorem[auto counter, number within=section, crefname={definition}{definitionen},
Crefname={Definition}{Definitionen}]{definition}{Definition}{
	breakable, colframe=green!70!black, fonttitle=\bfseries, coltitle = white, colbacktitle = green!70!black, enhanced, attach boxed title to top left={yshift=-0.1in, xshift=0.15in},}{def}

\newtcbtheorem[auto counter, number within=section, crefname={Satz}{Sätze},
Crefname={Satz}{Sätze}]{satz}{Satz}{%
	breakable, lower separated=false,
	colframe=red!70!white, fonttitle=\bfseries, coltitle = white, colbacktitle = red!70!white, enhanced, attach boxed title to top left={yshift=-0.1in,xshift=0.15in},
}{satz}




%%%%%%%%%%%%%%%%%%%%%%
%%%% Mathe-Befehle %%%
%%%%%%%%%%%%%%%%%%%%%%


\newcommand{\N}{\mathbb{N}}
\newcommand{\Z}{\mathbb{Z}}
\newcommand{\Q}{\mathbb{Q}}
\newcommand{\R}{\mathbb{R}}
\newcommand{\C}{\mathbb{C}}
\newcommand{\LL}{\mathbb{L}}
\newcommand{\DD}{\mathbb{D}}
\newcommand{\Ra}{\Rightarrow}

% TikZ
\usepackage{pgfplots}	
\pgfplotsset{compat=1.8}
\usepackage{pgf,tikz}
\usepackage{mathrsfs}
\usetikzlibrary{arrows}
\usetikzlibrary{trees}


\usetikzlibrary{arrows,positioning,shapes,fit,calc}

\pgfdeclarelayer{background}
\pgfsetlayers{background,main}
